\begin{pastexam}[label={px:ncu-114-dsa-10}]{中央資工所 114 學年度 第 10 題}
（複選；每題 5 分，全部答對才給分，答錯倒扣 1 分）

（第 8--10 題共用下面的說明和程式。）

Consider a UAV (Unmanned Aerial Vehicle) with F unit of fuel travels from NCU to a destination
``target'' km away. There are n gas stations along the way. When the UAV refuels at a gas station, all
the fuel of the gas station is transferred into the UAV.

Suppose that the UAV consumes one unit of fuel for every kilometer it travels. Let the position and fuel
of a gas station indicate the distance between the gas station and NCU and its volume of fuel,
respectively. Return the minimum number of refueling stops the UAV must make in order to reach its
destination.
\begin{minted}[fontsize=\small]{c}
typedef struct {
    int position;
    int fuel;
} stop_info;

int Refuel(int target, int F, int n, stop_info s[]) {
    int N_R = 0;   // number of refuel
    int i;   //gas stations ID
    stop_info X;

    while (F < target) {
        for (i=0;  L1  ; ++i)
            Q.push(s[i].fuel);
        if (Q.empty()) return -1;
        X=Q.pop();
        F += X.fuel
        N_R++;
    }
    return N_R;
}
\end{minted}

Consider Refuel algorithm above. Which of the following statements is true?
\begin{enumerate}[label=(\alph*)]
\item The time complexity of Refuel algorithm above is $\theta$(n$^2$logn)
\item The time complexity of the minimum number of refuel problem is $\theta$(n$^2$ logn)
\item The time complexity in Refuel algorithm above is $\theta$(nlogn)
\item The time complexity of the minimum number of refuel problem is $\theta$(nlogn)
\end{enumerate}
\end{pastexam}

\begin{pastsoln}
答案：(c)。

程式照歷屆試題~\ref{px:ncu-114-dsa-08} 詳解的讀法：加油站依位置排好，\texttt{i} 只在開始時設成 $0$，
\texttt{Q} 是最大堆積。 時間複雜度指最壞情況。

\textbf{Refuel 演算法是 $\Theta(n \log n)$}。 用整體來數，不要一輪一輪乘（資結書第~1.7~節）：
\begin{itemize}
\item \texttt{i} 從 $0$ 只往上加到最多 $n$，內層 \texttt{for} 在整個執行過程中總共只跑 $O(n)$ 次，
每個站最多推入一次：$n$ 次 \texttt{push}，每次 $O(\log n)$。
\item 每一輪 \texttt{while} 彈出一個站。 站最多 $n$ 個，所以 \texttt{while} 最多 $n + 1$ 輪（最後一輪
可能因為 \texttt{Q} 空了而回傳 $-1$），\texttt{pop} 最多 $n$ 次，每次 $O(\log n)$。
\end{itemize}
合計 $O(n \log n)$。 最壞情況也真的要這麼多：起始油量就到得了所有的站、站的油量由小到大排，而且
油量總和剛好在最後才夠，例如 $F = n$、$\mathit{target} = F + 1 + 2 + \cdots + n$、第 $i$ 站位置 $i$、
油量 $i$。 推入時每個新值都是最大的，要一路往上調到根，第 $j$ 次推入要調 $\lfloor\lg j\rfloor$ 層，合計
$\Theta(n \log n)$；而且 $n$ 個站都要用上，\texttt{pop} 也做 $n$ 次。 所以是 $\Theta(n \log n)$，
(c) 對。 若站沒有事先排序，先排序也只要 $\Theta(n \log n)$，結論不變。

\textbf{逐項}：
\begin{itemize}
\item (a) 把 \texttt{while}（最多 $n$ 輪）乘上內層 \texttt{for}（最多 $n$ 次）再乘 $\log n$，是把
\texttt{i} 誤當成每輪重來的上界，不是緊的界。 錯。
\item (b) 問題本身的複雜度指「解這個問題的最好演算法」的時間。 上面的演算法已經是 $O(n \log n)$，問題
不可能需要 $\Theta(n^2 \log n)$。 錯。
\item (c) 對，見上。
\item (d) 要說問題是 $\Theta(n \log n)$，除了有 $O(n \log n)$ 的演算法，還要證明\textbf{任何}演算法都
需要 $\Omega(n \log n)$。 這要另外做下界證明，題目和一般教材都沒有，上面的分析只說明演算法是 $\Theta(n \log n)$，不能拿來斷言問題本身的複雜度。 錯。
\end{itemize}
題目說「Which of the following statements is true?」，只有 (c)。
\end{pastsoln}
